% Generated draft: requires mathematical and source review.
\documentclass[11pt]{article}
\usepackage[margin=25mm]{geometry}
\usepackage{fontspec}
\setmainfont{Latin Modern Roman}
\usepackage{amsmath,amssymb,mathtools,unicode-math}
\setmathfont{Latin Modern Math}
\usepackage{enumitem,needspace,xurl,hyperref}
\setlength{\parindent}{0pt}
\setlength{\parskip}{5pt}
\setlength{\emergencystretch}{3em}
\newcommand{\N}{\mathbb{N}}
\newcommand{\Z}{\mathbb{Z}}
\newcommand{\Q}{\mathbb{Q}}
\newcommand{\R}{\mathbb{R}}
\newcommand{\C}{\mathbb{C}}
\newcommand{\F}{\mathbb{F}}
\newcommand{\A}{\mathbb{A}}
\newcommand{\PP}{\mathbb{P}}
\newcommand{\E}{\mathbb{E}}
\newcommand{\eps}{\varepsilon}
\newcommand{\dd}{\,\mathrm{d}}
\begin{document}
\section*{UnsolvedMath ID: 800650}
\textbf{Source problem code:} catalogue entry
\subsection{Definitions and mathematical statement}
The supplied excerpt does not specify the ambient class of spaces or define its cover notation. The following gives the conventional interpretation, subject to verification against the cited source. Let \(X\) be a topological space. A countable Borel cover of \(X\) is a countable family \(\mathcal U\) of Borel subsets of \(X\) such that \(\bigcup\mathcal U=X\) and \(X\notin\mathcal U\). Define the following classes of such covers:
\begin{itemize}
\item \(\mathcal B_\Lambda(X)\): large covers, meaning that each point of \(X\) belongs to infinitely many members of \(\mathcal U\).
\item \(\mathcal B_\Omega(X)\): \(\omega\)-covers, meaning that every finite subset of \(X\) is contained in some member of \(\mathcal U\).
\item \(\mathcal B_\Gamma(X)\): infinite \(\gamma\)-covers, meaning that each point of \(X\) belongs to all but finitely many members of \(\mathcal U\).
\item \(\mathcal B_{\mathrm T}(X)\): \(\tau\)-covers, meaning large covers for which, for every \(x,y\in X\), at least one of the sets
\[
\{U\in\mathcal U:x\in U,\ y\notin U\},\qquad
\{U\in\mathcal U:y\in U,\ x\notin U\}
\]
is finite.
\end{itemize}
For cover classes \(\mathcal A(X)\) and \(\mathcal D(X)\), use the subcover formulation of \(\mathsf{Split}(\mathcal A,\mathcal D)\): every \(\mathcal U\in\mathcal A(X)\) admits disjoint subfamilies \(\mathcal V_0,\mathcal V_1\subseteq\mathcal U\) with \(\mathcal V_0,\mathcal V_1\in\mathcal D(X)\). Equivalently, \(\mathcal U\) can be partitioned into two subfamilies, each containing a cover in \(\mathcal D(X)\). This convention also requires source verification.

A property \(P\) is hereditary on an ambient class \(\mathscr C\) if
\[
\forall X\in\mathscr C\ \forall Y\subseteq X,
\qquad P(X)\Longrightarrow P(Y),
\]
where \(Y\) carries the subspace topology and its Borel covers are defined relative to that topology. Here \(\mathscr C\) is a placeholder for the intended ambient class, which is not specified in the supplied record; no restriction to Borel subsets \(Y\) is imposed by this interpretation of hereditary.

\textbf{Problem (the catalogue statement).} Determine, separately for each of
\[
\mathsf{Split}(\mathcal B_\Omega,\mathcal B_\Lambda),\qquad
\mathsf{Split}(\mathcal B_\Omega,\mathcal B_\Omega),\qquad
\mathsf{Split}(\mathcal B_{\mathrm T},\mathcal B_{\mathrm T}),\qquad
\mathsf{Split}(\mathcal B_{\mathrm T},\mathcal B_\Gamma),
\]
whether it is hereditary in the intended ambient class. The missing ambient class and unverified notation conventions prevent this draft from specifying a unique unconditional formulation. The registry labels the problem open; that label is supplied source data, not an independently verified status assessment.
\subsection{Short English statement}
For each of four Borel cover splitting properties, ask whether every subspace of a space having the property also has it. The supplied excerpt leaves the intended class of spaces and the precise cover and splitting conventions unstated, so the expanded formulation requires source verification.
\subsection{Sources}
\begin{itemize}
\item Local UnsolvedMath registry, permanent ID \(800650\), exact source problem code \texttt{catalogue entry}; supplied complete source record, upstream revision \texttt{d371ee917ac9647fbb06cd6fc529deec6f3dda82}. The registry attributes the problem to Tsaban, gives proposed year 2003, and labels it open.
\item Tsaban, \emph{Selection Principles in Mathematics: A Milestone of Open Problems} (2003), Problem 23. Supplied source URL: \url{https://arxiv.org/abs/math/0312182}. The record reports access on 2026-10-16 and extraction by \texttt{arxiv\_latex}; these are supplied provenance claims, not an assertion of independent access or verification for this draft.
\item Supplied attribution/context: \texttt{[splittability, Problem 7.5]}. The record does not supply the bibliographic entry corresponding to \texttt{splittability}.
\end{itemize}
\end{document}
